\section{Complete CALM: Proofs and Details}
\label{app:proofs-sec3}

This appendix proves the results of Section~\ref{sec:complete-calm} and gives the formal
definition of the lifted output events of Section~\ref{sec:lifted-output-events}.

\subsection{Proofs}

\begin{proof}[Proof of Proposition~\ref{prop:common-expansion-form}]
First assume Lifted Monotonicity, and let $(i_1,o_1),\ldots,(i_k,o_k)$ be
rows with common \indesc $j$.  If $k=0$, any $u\in\fiber{j}$ will do, and
$\fiber{j}$ is nonempty since $\Spec$ is total.  Otherwise, since
$i_r\InputOrd j$, Relational Monotonicity gives, for each $r$, a
$u_r\in\fiber{j}$ with $o_r\OutputOrd u_r$.  The $u_r$ all lie in the one
fiber $\fiber{j}$, which is directed, so some $u\in\fiber{j}$ \outrel every
$u_r$, hence every $o_r$.  Now assume the common-extension condition, and let
$i\InputOrd j$ and $X\subseteq\fiber{i}$ be nonempty and finite.  The rows $\{(i,x):x\in X\}$
have common \indesc $j$, so some $u\in\fiber{j}$ \outrel all of $X$; that is
$\fiber{i}\LiftedOrd\fiber{j}$.
\end{proof}

The Survivor test (Proposition~\ref{prop:confluence}) and
Lemma~\ref{lem:run-residual} combine into the form that the later arguments
use.

\begin{lemma}[Reconciliation]
\label{lem:reconciliation}
Let runs end at $(i_r,R_r)$, $r=1,\ldots,m$, let $j$ be a common \indesc of
the $i_r$, and let $X$ be the set of outputs emitted on any of them.  The runs
are confluent at $j$ if and only if some $u\in\fiber{j}$ \outrel every output
in $X$.
\end{lemma}

\begin{proof}
By Lemma~\ref{lem:run-residual}, $(j,u)\in R_r$ exactly when $u\in\fiber{j}$
\outrel every output emitted on run $r$, and Proposition~\ref{prop:confluence}
gives the claim.
\end{proof}

\begin{proof}[Proof of Proposition~\ref{prop:exposure-boundary}]
By Proposition~\ref{prop:common-expansion-form} it suffices to treat the
common-extension condition.  First assume it holds.  Let $(i,R)$ be reached by
a run (\commitment-free or not), and let $o\in R(i)$, $j\InputGe i$ legal, and
$x\in R(j)$.  The rows $(i,o)$ and $(j,x)$ of $P$ have common \indesc $j$, so
the condition supplies $u\in P(j)$ \outrelpart both $o$ and $x$.  By
Lemma~\ref{lem:run-residual}, $(j,u)\in R$, since $j\InputGe i$ and $u$
\outrel $o$, which \outrel every output emitted on the run.  So $u\in R(j)$
\outrel $o$ and $x$, and the legal-output test
(Proposition~\ref{prop:exposability}) makes $o$ a legal output.  Now assume
every output at every state reachable without \commitment is legal, and let
rows
$(i_1,o_1),\ldots,(i_k,o_k)$ have common \indesc $j$.  By totality there is some
$u_0\in P(j)$.  For $r=1,\ldots,k$ in turn, the legal-output test for $o_r$ at
$\sigma_P(i_r)$, applied with $x=u_{r-1}$, gives $u_r\in P(j)$ \outrelpart both
$o_r$ and $u_{r-1}$.  Then $u_k$ \outrel every $o_r$, and the
common-extension condition follows.
\end{proof}

\begin{proof}[Proof of Theorem~\ref{thm:complete-calm}]
The proof of Lemma~\ref{lem:confluence-legality} covers nondeterministic specs,
so the runs are consistent exactly when no output event of a run applies a
\commitment, that is, when every output that a run can emit is legal.  As in
the proof of Theorem~\ref{thm:functional-calm}, this holds exactly when every output at every state reachable without \commitment is legal.  By
Proposition~\ref{prop:exposure-boundary}, that holds exactly when the spec
satisfies Lifted Monotonicity.
\end{proof}

\paragraph{Pairs suffice for the common-extension condition.}
\label{app:pairs-suffice}
Suppose every two rows with a common \indesc $j$ have a common \outnoun in
$\fiber{j}$.  Then the condition holds for every finite family of rows with
common \indesc $j$.  We replace the first two rows of the family by a row
$(j,u)$ whose output $u$ \outrel both, and repeat until one row remains; a
single row, paired with any row at $j$, has an \outnoun in $\fiber{j}$.  The
pairwise hypothesis must cover all rows of $P$, including rows at $j$ itself,
since the replacement rows are at $j$.  Checking only rows at single-proposal
inputs does not suffice, even with families of up to $k$ rows.  In the
$k$-set agreement example of Appendix~\ref{app:additional-applications},
every family of at most $k$ rows at such inputs is compatible at each common
\indesc, but the $k+1$ proposal rows are not.

\subsection{Lifted output events}
\label{app:lifted-output-events}

Fix a state $(i,R)$ and an output $o\in R(i)$.  A row is \emph{ruled out by}
$o$ when its output does not \outrelbase $o$.  The ruled-out rows fall into two
disjoint sets:
\begin{align}
  \Subsumed_o&=\{(k,x)\in R\setminus R[\mathrm{out}\OutputGe o]:\notag\\
    &\qquad\quad x\OutputOrd u\text{ for some }u\in R[\mathrm{out}\OutputGe o](k)\},
      \label{eq:output-subsumptions}\\
  \Excluded_o&=(R\setminus R[\mathrm{out}\OutputGe o])\setminus\Subsumed_o.
      \label{eq:commitment-removals}
\end{align}

\begin{definition}[Lifted output event, output subsumption]
\label{def:lifted-output-event}
Let $(i,R)$ be a state and $o\in R(i)$.  The \emph{output event}
$\mathsf{output}(o)$ keeps exactly the rows whose outputs \outrelbase $o$ and
emits $o$, atomically, ending at $(i,R[\mathrm{out}\OutputGe o])$ as in
Equation~\ref{eq:output-action}.  It removes each row $(k,x)$ in
$\Subsumed_o$, for which some $u\in R(k)$ \outrel both $x$ and $o$, by
\emph{output subsumption}, which is not a \commitment.
\end{definition}

\begin{definition}[Canonical \commitment]
\label{def:canonical-commitment}
When $\Excluded_o\neq\varnothing$, the output event applies the \emph{canonical
\commitment} $x_o$:
\begin{equation}
  x_o(i,R)=\bigl(i,\ R\setminus\Excluded_o\bigr).
  \label{eq:canonical-commitment-relational}
\end{equation}
  Thus $x_o$ removes exactly the rows in $\Excluded_o$.  At an input $k$, if
  $\Excluded_o(k)=R(k)$, it rules out $k$; otherwise it rules out only the
  outputs in $\Excluded_o(k)$.  A row in $\Subsumed_o(k)$ has a kept \outnoun
  at $k$, so removing it by output subsumption does not rule out $k$.
\end{definition}

In Figure~\ref{fig:ambiguity-frames}(f), emitting $[A,B]$ places the row for
$[B,A]$ in $\Excluded_o$, so the event rules out that output at $\{A,B\}$.
In panel~(g), the removed route sets, at both inputs, are in $\Subsumed_o$,
each \outrelpast by a retained route set at its input, so the event subsumes them without
\commitment.  With one row per input, $\Subsumed_o$ is empty, and $x_o$ rules
out each input whose row is removed; this is the functional output event of
Definition~\ref{def:output-event}.  Apart from output subsumption and
\commitments, a row leaves the remainder table only by input loss
(\S\ref{sec:functional-events}).

Legality lifts with the event: $o$ is a legal output when $\Excluded_o$ is
empty, and then $\mathsf{output}(o)$ applies no \commitment, though it may
still subsume the rows in $\Subsumed_o$.  The legal-output test lifts as well:

\begin{proposition}[Legal-output test]
\label{prop:exposability}
An output $o\in R(i)$ is legal at $(i,R)$ if and only if, for every legal
$j\InputGe i$ and $x\in R(j)$, some $u\in R(j)$ \outrel both:
\begin{equation}
  o\OutputOrd u
  \qquad\text{and}\qquad
  x\OutputOrd u.
  \label{eq:safe-exposure}
\end{equation}
\end{proposition}

\begin{proof}
The output $o$ is legal exactly when $\Excluded_o$ is empty, that is, when
every ruled-out row $(j,x)$ is \outrelpast by some row $(j,u)$ \outrelpart
$o$.  Rows that already \outrelbase $o$ satisfy the test with $u=x$, so this
is the stated test.
\end{proof}
